\begin{abstract}
Predicting the remaining useful life (RUL) of commercial turbofan engines from in-flight sensor recordings is a core problem in aerospace prognostics: an accurate, well-calibrated estimator enables condition-based maintenance, reduces unscheduled removals, and bounds the operational risk of in-service failure. The N-CMAPSS dataset~\cite{ariaschao2021ncmapss} replaces the synthetic CMAPSS benchmark~\cite{saxena2008damage} with run-to-failure trajectories generated under real recorded flight envelopes, yielding variable-length flights of $10^4$--$10^5$ timesteps per engine and heterogeneous operating conditions that previously published deep models were not designed to ingest. We present a deep-learning pipeline for variable-length N-CMAPSS flights and use it to compare recurrent, temporal-convolutional, transformer, and multi-scale attention families under the same preprocessing, training, and evaluation code. Operationally short windows of roughly one thousand timesteps are competitive with, and in our benchmarks substantially better than, full-flight inputs; the gap is confounded with architecture, but the pattern is consistent with the short-window choices of prior PHM-2021 challenge submissions~\cite{solismartin2021stacked,lovberg2021variable}. A multi-scale TCN with global-fusion attention (MSTCN) sits in the top performance cluster alongside transformer and WaveNet baselines, is best by RMSE on FD2, and outperforms full-flight RNN baselines whose $R^2$ is negative. We use asymmetric MSE with a 2$\times$ late-prediction penalty as a design choice anchored on the asymmetric PHM scoring function~\cite{saxena2008damage} and recent loss-design work~\cite{mpm2024loss}, rather than as an empirical contribution. On the apples-to-apples FD1 benchmark, the best controlled model attains RMSE 6.523 cycles, within the band reported by N-CMAPSS prior art.
\end{abstract}

\section{Introduction}

Modern commercial turbofans are instrumented with dense sensor suites whose readings, integrated over thousands of flight hours, encode the gradual degradation of compressor, turbine, and combustor components. Translating those recordings into a calibrated estimate of remaining useful life is the central task of data-driven engine prognostics: a reliable RUL estimate lets airlines plan maintenance against actual wear rather than fixed intervals, lowers the rate of unscheduled removals, and bounds the operational risk of late detection of an evolving fault. The cost structure of this task is sharply asymmetric. Predicting RUL too early is conservative and incurs maintenance overhead; predicting it too late risks an in-service failure. The PHM-2008 scoring function~\cite{saxena2008damage} encodes this asymmetry directly and remains the canonical evaluation metric in the field.

For more than a decade, deep-learning research on this problem has been driven by the synthetic CMAPSS benchmark~\cite{saxena2008damage}, where each engine produces a few hundred cycles of low-dimensional sensor data under a small set of operating conditions. The N-CMAPSS dataset of Arias-Chao et al.~\cite{ariaschao2021ncmapss} was introduced to close the gap between this benchmark and the conditions encountered in commercial operation: trajectories are simulated from real recorded flight envelopes spanning ascent, cruise, and descent, with sample-rate-level sensor data on the order of $10^4$--$10^5$ timesteps per cycle. The hybrid physics-augmented baseline of Arias-Chao et al.~\cite{ariaschao2022fusing} and the published submissions to the 2021 PHM Society Data Challenge~\cite{lovberg2021variable,devol2021inception,solismartin2021stacked} establish N-CMAPSS as the contemporary benchmark of record for turbofan RUL.

N-CMAPSS poses three difficulties that did not arise on CMAPSS. Flights are variable-length, with each engine yielding a different total cycle count and per-cycle sequence length. Operating conditions are heterogeneous within and across flights, so a single model must generalize across the full envelope. And the sequence regime is qualitatively different: at $10^4$--$10^5$ timesteps per cycle, naively ingesting full flights pushes recurrent and quadratic-attention models toward their failure modes for long inputs.

The methodological response to these difficulties has spanned four families. Recurrent networks such as LSTM~\cite{hochreiter1997lstm}, GRU~\cite{cho2014gru}, and bidirectional variants were the first applied. Temporal convolutional networks~\cite{bai2018tcn}, WaveNet-style dilated convolutions~\cite{oord2016wavenet}, and dilated/inception variants~\cite{lovberg2021variable,devol2021inception} were introduced to handle long inputs without sequential bottlenecks. Self-attention models, building on the Transformer~\cite{vaswani2017attention} and long-sequence variants such as Informer~\cite{zhou2021informer}, have been adapted to RUL through hierarchical and two-stage designs~\cite{fan2024star,ttsnet2025}. Most recently, multi-scale TCN architectures with attention modules over the dilation hierarchy~\cite{xu2024msattn} have been proposed as a way to combine the long effective receptive field of TCNs with the cross-channel and cross-scale weighting of attention. The publicly reported benchmarks across these families are not directly comparable: most are reported on CMAPSS rather than N-CMAPSS, training protocols and input lengths vary substantially, and few studies report the same accuracy, RMSE, and PHM-score triple under a single harness.

This paper contributes a deep-learning pipeline for variable-length N-CMAPSS flights together with a shared registry of 22 architectures drawn from the main recurrent, convolutional, hybrid, transformer, and multi-scale attention families. The pipeline standardizes sequence preparation (last-$N$ truncation, per-feature standardization), the loss function (asymmetric MSE with a 2$\times$ late-prediction penalty), the evaluation metrics (RMSE, $R^2$, PHM score, accuracy at $\pm 20$ cycles), and the apples-to-apples experimental harness used for the head-to-head numbers. The results support the following claims:

\begin{itemize}
  \item \textbf{C1 (sequence length).} Operationally short windows of roughly one thousand timesteps are competitive with, and in our benchmarks substantially better than, full-flight inputs. Full-flight RNN baselines yield negative $R^2$ (T8: LSTM RMSE 22.28, $R^2 = -0.066$; T20: LSTM RMSE 22.28; T21: BiLSTM RMSE 22.20) while short-window models cluster at RMSE 6.5--7.6 (T1--T3); the magnitude is confounded with architecture choice and we report it as a benchmark ceiling rather than as a controlled single-model sweep.
  \item \textbf{C2 (multi-scale temporal attention).} A multi-scale TCN with global-fusion attention (MSTCN) sits in the top cluster of attention and multi-scale models on N-CMAPSS, is best on FD2 by RMSE (T11: 6.495), and is far ahead of the full-flight RNN baselines. We do not claim MSTCN beats every alternative: in the apples-to-apples FD1 benchmark MSTCN ranks third (T3: 7.604) behind WaveNet (T1: 6.523) and CNN-GRU (T2: 7.006).
  \item \textbf{C3 (asymmetric loss).} We adopt asymmetric MSE with a 2$\times$ late-prediction penalty as a methodological choice anchored on the PHM scoring asymmetry of Saxena et al.~\cite{saxena2008damage} and the recent multiple-penalty-mechanism loss of \cite{mpm2024loss}, not as an empirical contribution. A controlled symmetric-MSE baseline and an $\alpha$-sensitivity sweep remain to be run.
\end{itemize}

The rest of the paper is organized as follows. Section~2 reviews related work in N-CMAPSS prognostics, attention-based RUL models, and asymmetric loss design. Section~3 describes the pipeline, including data handling, the asymmetric loss, and the MSTCN architecture. Section~4 presents the apples-to-apples benchmark across the registered architecture set. Section~5 discusses C1--C3 and the missing ablations. Section~6 concludes.
